\section{Resultados}

\subsection{Fase 1: documentos exigidos en los trámites y su clasificación}

Los resultados de esta fase son preliminares: la codificación la realizó un solo investigador y la atribución del emisor de varios documentos debe confirmarse con la norma de cada trámite \pendiente{DATO PENDIENTE: confirmación del emisor de los documentos marcados y, si se decide, revisión por un segundo revisor}.

\subsubsection{Inventario de trámites}

El buscador del Portal Único del Estado Colombiano registraba, el 31 de octubre de 2026, 85 trámites a nombre de la Gobernación del Huila \parencite{govco2026fichas}, los mismos que muestra el Visor del SUIT de Función Pública \parencite{funcionpublica2026visorsuit}. Al emparejarlos con el catálogo de la sede electrónica de la entidad \parencite{gobernaciondelhuila2026tramites} se encontró lo siguiente:

\begin{itemize}
\item Siete trámites registrados en el SUIT no aparecen en el catálogo de la sede electrónica: la pensión de jubilación y las cesantías parciales para docentes oficiales, la asistencia técnica a establecimientos educativos, la clasificación en el régimen de educación de un establecimiento privado, la cancelación de la personería jurídica de ligas y clubes deportivos, las facilidades de pago para deudores de obligaciones no tributarias y la autorización de funcionamiento de establecimientos farmacéuticos.
\item Un trámite del catálogo, la inscripción de dignatarios de las organizaciones comunales, no figura en el SUIT a nombre de la Gobernación: la sede electrónica lo enlaza a una ficha de gov.co que responde que el trámite no existe. En el extracto del SUIT de 2022 sí figuraba, con tres requisitos documentales, todos de la categoría (c) \parencite{funcionpublica2022suit}.
\item Cinco trámites del catálogo no tenían código y otro tenía el código de un trámite distinto; sus códigos se identificaron por el nombre.
\end{itemize}

Las fuentes no coinciden en la modalidad. El catálogo de la sede electrónica clasificaba sus 80 registros como 27 en línea, 12 parcialmente en línea y 41 presenciales; el SUIT clasifica los 85 trámites como 29 totalmente en línea, 24 parcialmente en línea y 32 sin medio electrónico disponible.

El SUIT registra, además, la fecha de la última modificación de cada trámite. De los 85, 26 se modificaron en 2026, 8 en 2025 y 16 entre 2022 y 2024; los 35 restantes no se modifican desde 2021 o antes, y 14 de ellos tienen su última modificación en 2015.

\subsubsection{Requisitos registrados}

Los registros de los 85 trámites en el SUIT contienen 760 requisitos: 472 de tipo documento, 144 verificaciones que realiza la propia entidad, 124 canales de solicitud, 10 pagos y 10 formularios. El propio SUIT clasifica 171 de los 472 documentos como estandarizados y 301 como propios del trámite.

Al aplicar las reglas de codificación, los 472 registros de tipo documento se redujeron a 324 requisitos documentales distintos en 64 trámites. La reducción se debe a que el registro repite un mismo documento para distintas personas (por ejemplo, la cédula del educador fallecido y la del beneficiario), registra por separado documentos alternativos (cédula de ciudadanía o de extranjería) y registra como documento requisitos que no lo son: se excluyeron del conteo 30 solicitudes, 10 recibos o declaraciones de pago, un formato de la entidad y otros dos registros que no son documentos.

Veintiún trámites no exigen ningún documento distinto de la solicitud o del pago. En diez no hay ningún documento registrado; en ocho, el único documento registrado es la carta o el formato de solicitud, y en tres, el recibo o la declaración del pago. Entre ellos están el cambio de sede y la clausura de un establecimiento educativo, y la certificación de prestación de servicios para contratistas, que el Programa de Transparencia y Ética Pública 2026 prioriza para racionalización \parencite[p.~135]{gobernaciondelhuila2026ptep}. La sede electrónica no permite resolver esa duda para 20 de los 21 trámites, porque los enlaza a su ficha de gov.co, que reproduce el registro del SUIT. La excepción es la inscripción y concepto sanitario de establecimientos farmacéuticos, que tiene una página propia en la sede electrónica: allí se indica que, para varias condiciones, el solicitante «no debe presentar ningún soporte ya que la institución verificará que usted la cumpla», y se exigen, entre otros, el certificado de uso del suelo y el certificado de idoneidad sanitaria, que no figuran en el registro del SUIT \parencite{gobernaciondelhuila2026tramites}. Además, siete de los diez trámites sin ningún documento en 2026 tampoco tenían documentos registrados en el extracto de 2022, y los otros tres no existían entonces.

\subsubsection{Clasificación según quién expide el documento}

La Tabla~\ref{tab:categorias} presenta la clasificación de los 324 requisitos documentales. Algo más de la mitad corresponde a documentos que expide otra autoridad pública o un registro público, y algo menos de uno de cada cinco, a documentos que expide o custodia la propia Gobernación.

\begin{table}[htbp]
\centering
\caption{Requisitos documentales según quién expide o custodia el documento}
\label{tab:categorias}
\begin{tabular}{p{5.6cm}rrrrrr}
\hline
Categoría & \multicolumn{2}{c}{Todos los requisitos} & \multicolumn{2}{c}{Sin emisor variable} & \multicolumn{2}{c}{Criterio estricto} \\
 & n & \% & n & \% & n & \% \\
\hline
(a) Los expide o custodia la Gobernación & 56 & 17,3 & 23 & 10,0 & 2 & 1,0 \\
(b) Los expide otra autoridad o un registro público & 167 & 51,5 & 115 & 50,2 & 115 & 55,3 \\
(c) Solo los puede aportar el ciudadano & 101 & 31,2 & 91 & 39,7 & 91 & 43,8 \\
\hline
Total & 324 & 100,0 & 229 & 100,0 & 208 & 100,0 \\
\hline
\end{tabular}

\smallskip
\footnotesize\textit{Nota.} «Sin emisor variable» excluye los 95 requisitos cuyo emisor depende del caso. «Criterio estricto» excluye, además, los 21 requisitos que se presumen expedidos por la Gobernación por ser actos previos de la misma competencia del trámite. Elaboración propia con datos del Visor del SUIT \parencite{funcionpublica2026visorsuit}, consultado el 31 de octubre de 2026.
\end{table}

Dentro de la categoría (b), 74 requisitos corresponden a registros públicos, como el documento de identidad, el registro civil, el certificado de tradición y libertad, la licencia de tránsito o el RUT, y 15 los expiden particulares que ejercen funciones públicas, como las cámaras de comercio, las notarías o los centros de diagnóstico automotor. \pendiente{CONFIRMAR: según el PETI 2026, la Gobernación también expide certificados de existencia y representación legal de entidades sin ánimo de lucro y organizaciones comunales; revisar si en algún trámite ese certificado lo expide la propia entidad} Dentro de la categoría (c), 38 son documentos de terceros privados y 16 son declaraciones del ciudadano o de terceros. Del total, 136 requisitos son condicionales: solo se exigen en ciertos casos, según el registro. La proporción de cada categoría es similar entre los requisitos condicionales y los demás: 19,1\,\%, 50,7\,\% y 30,1\,\% en los condicionales, frente a 16,0\,\%, 52,1\,\% y 31,9\,\% en los demás; la condicionalidad se identificó por expresiones del texto del registro, como «si», «en caso de» o «cuando», o por la excepción que registra el SUIT. No se marcó ningún documento como expedido por una entidad descentralizada del departamento; el emisor del reconocimiento deportivo está por confirmar \pendiente{CONFIRMAR: emisor del reconocimiento deportivo, trámite T23048}. La distinción entre entidades integradas o no al servicio ciudadano digital de interoperabilidad queda pendiente de la respuesta de la Gobernación a la solicitud de información.

La clasificación es sensible a dos supuestos. El primero son los 95 requisitos cuyo emisor depende del caso, como el certificado de tiempo de servicio, el certificado de salarios o el diploma: si se excluyen, la proporción de la categoría (a) baja del 17,3\,\% al 10,0\,\%. El segundo es la presunción de que los actos previos de la misma competencia del trámite, como la resolución que reconoció la pensión de un docente o el acto de su retiro, los expidió la Gobernación: 21 de los 23 requisitos que quedan en la categoría (a) dependen de ella, y si también se excluyen, solo quedan 2 requisitos en esa categoría, el certificado de habilitación y la resolución de autorización de etiquetas. En los tres escenarios, la suma de las categorías (a) y (b), que agrupa los documentos que expide o custodia una autoridad pública o un registro público, queda entre el 56\,\% y el 69\,\%.

El SUIT indica el tipo de emisor en 71 de los 472 registros de documentos. La clasificación del estudio coincide con los 27 que el SUIT marca como de emisor público; en 33 el SUIT indica que el emisor puede ser público o privado, lo que concuerda con la marca de emisor variable, y en 3 registros del certificado de ingresos y retenciones el SUIT indica al ciudadano como emisor, mientras que el estudio lo clasificó en la categoría (b) por expedirlo el pagador de la prestación.

De los 64 trámites con requisitos documentales, 25 exigen al menos un documento de la categoría (a), 56 al menos uno de las categorías (a) o (b), y 8 solo exigen documentos de la categoría (c). Cada trámite exige entre 1 y 17 documentos, con una mediana de 4. Con el criterio estricto, solo 2 trámites conservan documentos de la categoría (a).

\subsubsection{Documentos que más se repiten}

La Tabla~\ref{tab:repetidos} muestra los documentos exigidos en más trámites, que son el insumo del segundo criterio de priorización.

\begin{table}[htbp]
\centering
\caption{Documentos exigidos en siete o más trámites}
\label{tab:repetidos}
\begin{tabular}{p{8.5cm}cr}
\hline
Documento & Categoría & Trámites \\
\hline
Documento de identidad & (b) & 30 \\
Certificado de existencia y representación legal & (b) & 11 \\
Certificado de tiempo de servicio & (a)\textsuperscript{1} & 10 \\
Diploma o acta de grado & (b)\textsuperscript{1} & 10 \\
Registro civil de nacimiento & (b) & 9 \\
Declaraciones de terceros o manifestaciones del solicitante & (c) & 9 \\
Resolución de pensión del docente & (a)\textsuperscript{2} & 8 \\
Certificado de salarios & (a)\textsuperscript{1} & 8 \\
Acto administrativo de retiro del servicio docente & (a)\textsuperscript{2} & 7 \\
Tarjeta profesional & (b)\textsuperscript{1} & 7 \\
Certificado de administradora de pensiones & (b)\textsuperscript{1} & 7 \\
Contrato & (c) & 7 \\
\hline
\end{tabular}

\smallskip
\footnotesize\textit{Nota.} Trámites: número de trámites, de los 64 con requisitos documentales, que exigen el documento. \textsuperscript{1}~Emisor variable. \textsuperscript{2}~Se presume expedido por la Gobernación por tratarse de un acto previo de la misma competencia del trámite; por confirmar con la norma del trámite. Elaboración propia con datos del Visor del SUIT \parencite{funcionpublica2026visorsuit}, consultado el 31 de octubre de 2026.
\end{table}

Los documentos de la categoría (a) se concentran en los trámites de prestaciones de los docentes oficiales: 40 de los 56 corresponden a los 11 trámites de pensiones, cesantías, reliquidación, auxilio funerario y seguro por muerte. En esos trámites se pide al solicitante, entre otros, la resolución que reconoció la pensión, el acto de retiro del servicio, el certificado de salarios y el certificado de tiempo de servicio, documentos que se presumen expedidos por la Gobernación \pendiente{CONFIRMAR: emisor de estos actos según la norma de cada trámite}. Esta concentración depende de los dos supuestos anteriores: sin emisor variable, 17 de los 23 documentos de la categoría (a) están en esos trámites, y todos ellos dependen de la presunción de la misma competencia.

\subsubsection{Verificaciones que realiza la propia entidad}

El SUIT registra 144 verificaciones que la Gobernación realiza por su cuenta, en 48 trámites. Su distribución es desigual: 23 trámites tienen una sola verificación y uno, el ascenso en el escalafón nacional docente, concentra 43. En seis trámites la entidad registra como verificación que el solicitante esté inscrito en la cámara de comercio y, en cinco de ellos, no exige el certificado de existencia y representación legal. En diez trámites de prestaciones docentes registra como verificación el tiempo de servicio y, en ocho de ellos, exige a la vez el certificado de tiempo de servicio.

\subsubsection{Comparación con el SUIT de 2022}

El extracto del SUIT publicado en datos abiertos, con datos de febrero de 2022, registraba 87 trámites de la Gobernación \parencite{funcionpublica2022suit}. Entre esa fecha y octubre de 2026 salieron del registro ocho trámites, cinco de ellos de organizaciones comunales, y entraron seis: las cesantías parciales para docentes, las facilidades de pago de obligaciones no tributarias, la asistencia técnica a establecimientos educativos, la clasificación en el régimen de educación, la cancelación de la personería jurídica de ligas y clubes deportivos, y la inscripción y concepto sanitario de establecimientos farmacéuticos.

Para los 79 trámites presentes en ambas fechas se compararon los requisitos documentales codificados con las mismas reglas. En 70 trámites el conjunto de documentos exigidos es idéntico. En los otros nueve se añadieron 13 documentos y se retiraron 7; en tres de ellos la última modificación registrada es anterior a febrero de 2022, por lo que la diferencia puede deberse a la forma en que cada fuente publica los nombres y no a un cambio de requisitos. Los documentos de la categoría (a) en estos 79 trámites pasaron de 51 a 52. En conjunto, entre 2022 y 2026 la Gobernación no retiró de sus trámites ningún documento expedido por ella misma.

\subsection{Fase 2: relación de los documentos con las bases de datos y los mecanismos de intercambio}

Esta fase se aplicó a los 223 requisitos de las categorías (a) y (b), que corresponden a 69 documentos distintos. La relación entre cada documento y una base de datos se basa en lo que declaran el documento de datos maestros de referencia de la Gobernación \parencite{gobernaciondelhuila2026datosmaestros} y el catálogo de sistemas de información del PETI 2026 \parencite[pp.~34-37]{gobernaciondelhuila2026peti}, cuyo texto se obtuvo por reconocimiento óptico de caracteres. El estudio no verificó qué contienen realmente esos sistemas. Además, el documento de datos maestros describe las tablas de datos de la Gobernación y no sus trámites, de modo que lo que no menciona no puede leerse como ausencia: en adelante, «no declarado» significa que las fuentes no lo mencionan, no que no exista. Ese documento no indica fecha, autor ni dependencia responsable, y su numeración de tablas se repite, por lo que se cita por página \pendiente{CONFIRMAR: fecha, autoría y vigencia del documento de datos maestros; no se incluyó en la solicitud radicada y se prevé un seguimiento}. Para los documentos de otras autoridades se revisó, además, si la entidad que los expide ofrece en su portal una consulta pública en línea del mismo dato. La relación de cada requisito con su base, su mecanismo y su consulta pública se presenta en el Anexo~\pendiente{DATO PENDIENTE: número del anexo con la matriz de la Fase 2}.

\subsubsection{Documentos que expide la Gobernación}

De los 56 requisitos de la categoría (a), solo uno corresponde a un dato que una base declarada contiene por completo: el certificado de habilitación de prestadores de servicios de salud, cuyo estado y fecha registra la tabla de instituciones prestadoras de servicios de salud \parencite[p.~2]{gobernaciondelhuila2026datosmaestros}, aunque la fuente indica que ese dato lo toma del registro del Ministerio de Salud y Protección Social. En 43 requisitos, la base declarada contiene solo parte del dato o solo para una parte de los casos. En los 12 restantes, como la tornaguía, los reportes anuales de cesantías o la resolución de cesantías pagadas, las fuentes no declaran ninguna base que contenga el dato; en el caso de las cesantías, la tabla de historial laboral declara entre sus usos el cálculo de prestaciones sociales \parencite[p.~2]{gobernaciondelhuila2026datosmaestros}, pero no el reporte ni la resolución que piden los trámites. La Tabla~\ref{tab:cobertura} muestra estos resultados con los tres escenarios de la Fase 1: con todos los requisitos, sin los de emisor variable y con el criterio estricto, que además excluye los actos previos de la misma competencia del trámite. En el escenario estricto, la categoría (a) se reduce a 2 requisitos.

Los casos parciales se concentran en las prestaciones de los docentes. La tabla de docentes del sector educativo registra la fecha de ingreso y el grado en el escalafón, pero no el tiempo de servicio por entidad ni el salario; la tabla de servidores públicos sí registra el salario, pero no incluye a los docentes, y la tabla de pensionados registra el acto administrativo de reconocimiento de los pensionados a cargo del departamento \parencite[p.~2]{gobernaciondelhuila2026datosmaestros}. La tabla de historial laboral, que es de servidores, declara entre sus usos la \enquote{elaboración de certificaciones de tiempo de servicio y salarios para trámites de pensión} y describe así su mecanismo: \enquote{La entrega de certificaciones se hace por oficio o correo institucional} \parencite[p.~2]{gobernaciondelhuila2026datosmaestros}. El estudio supone que esa tabla respalda los certificados de los trámites docentes, pero la fuente no dice que cubra a los docentes. Lo que sí describe la fuente es un circuito interno en el que la dependencia que custodia el dato produce un certificado, y no una consulta del dato desde el trámite. En el documento de datos maestros, la expresión \enquote{cargue de archivo} aparece 14 veces como mecanismo; las tablas en que aparece corresponden a reportes de municipios, inventarios enviados a entidades nacionales o formatos de gestión interna, y ninguna se relacionó con un documento exigido en los trámites.

\subsubsection{Documentos que expiden otras autoridades}

De los 167 requisitos de la categoría (b), 39 (23,4\,\%) corresponden a datos de entidades con las que la fuente declara un servicio web: el RUNT, el RUP (que la fuente atribuye a \enquote{CCE}) y la DIAN \parencite[p.~1]{gobernaciondelhuila2026datosmaestros}; Colpensiones y el SIMAT del Ministerio de Educación \parencite[p.~2]{gobernaciondelhuila2026datosmaestros}, y el registro de talento humano en salud (RETHUS) y el IGAC, para los predios rurales \parencite[p.~3]{gobernaciondelhuila2026datosmaestros}. La fuente no declara el uso de una plataforma común ni dice que las conexiones sean bilaterales, y el estudio no pudo determinarlo. Solo en el caso de la DIAN la fuente describe un dato completo, el RUT, pero declara una validación con la API del RUT en procesos de contratación y de rentas, y no que se use en los trámites que lo exigen. En los demás casos el servicio no cubre el dato completo: el del RUNT no registra el número de la licencia de tránsito; el de Colpensiones sirve para validar derechos de los pensionados del departamento y no incluye a los fondos privados; el del RUP no equivale al certificado de existencia y representación legal; el de RETHUS solo cubre las profesiones de la salud, y el del IGAC solo los predios rurales.

En los otros 128 requisitos (76,6\,\%) las fuentes no declaran un mecanismo de intercambio. En 32 de ellos, la entidad que expide el documento ofrece una consulta pública en línea del mismo dato, pensada para una persona que consulta un documento a la vez y no para el intercambio entre sistemas. La Procuraduría General de la Nación indica que su certificado de antecedentes es gratuito, se obtiene de inmediato en su portal y que, desde la Ley 1238 de 2008, las entidades pueden consultarlo para quienes van a tomar posesión de un cargo o a contratar \parencite{procuraduria2026antecedentes}. La Contraloría General de la República expide en línea, sin costo, el certificado de antecedentes fiscales y afirma que \enquote{las entidades públicas deberán consultar directamente} el boletín de responsables fiscales cuando se trate de futuros contratistas o funcionarios \parencite{contraloria2026antecedentes}. Ambos certificados se exigen en un mismo trámite, el registro de sujetos pasivos del impuesto al consumo, que no es una posesión ni una contratación; el estudio no determinó si esa consulta puede usarse en ese trámite. Las otras 30 consultas corresponden al documento de identidad, el requisito más frecuente de la categoría (30 trámites). La única fuente es una noticia de 2013 de la Registraduría Nacional del Estado Civil, según la cual el certificado de vigencia de la cédula se obtiene gratis y por internet \parencite{registraduria2013vigencia}. El estudio no confirmó que ese servicio siga activo; en todo caso, el certificado confirma la vigencia y no reemplaza la copia del documento. Para los registros civiles de nacimiento, matrimonio y defunción (20 requisitos), la misma noticia solo ofrece a los ciudadanos consultar la oficina donde está inscrito el registro, pero declara que pone el Archivo Nacional de Identificación y el Sistema de Registro Civil a disposición de las demás entidades públicas, con fundamento en el artículo 227 de la Ley 1450 de 2011. Tampoco se pudo establecer si hoy se mantiene ni si la Gobernación tiene acceso, y el documento de datos maestros no menciona ninguna conexión con la Registraduría. Otras dos consultas complementan documentos que ya tienen un servicio web declarado: el RUES ofrece a los funcionarios de entidades información de los registros empresariales, sin costo, aunque no es el certificado de existencia y representación legal \parencite{rues2026consulta}, y el RETHUS permite consultar si el personal de salud cumple los requisitos para ejercer \parencite{minsalud2026rethus}.

Los 96 requisitos restantes (57,5\,\%) no tienen mecanismo declarado ni consulta pública identificada; los más frecuentes son el diploma o acta de grado (10 trámites) y los registros civiles (20 requisitos, entre ellos 9 de nacimiento). La Tabla~\ref{tab:cobertura} resume la situación del dato. Que el mecanismo no esté declarado no implica que no exista, y el estudio no puede determinar con fuentes públicas si falta.

La fuente niega en dos tablas una conexión con entidades de las que proceden documentos de esta categoría, pero en ninguno de los dos casos la negación se refiere a esos documentos. En la tabla de farmacias, la fuente indica que la habilitación \enquote{se notifica al INVIMA por oficio} y que no existe servicio web con esa entidad \parencite[p.~3]{gobernaciondelhuila2026datosmaestros}; la tabla trata de los establecimientos y no del registro sanitario ni del certificado de buenas prácticas de manufactura. En la tabla del catálogo de faltas disciplinarias, la fuente indica que no existe servicio web con la Procuraduría \parencite[p.~5]{gobernaciondelhuila2026datosmaestros}; ninguna tabla trata del certificado de antecedentes disciplinarios. Ninguna de las fuentes declara el uso de la plataforma X-Road ni del servicio ciudadano digital de interoperabilidad, y el documento de datos maestros indica en 46 ocasiones, en sus distintas tablas, que no existe servicio web, X-Road, enlace de base de datos ni conexión P2P.

\begin{table}[htbp]
\centering
\caption{Requisitos de las categorías (a) y (b) según la situación del dato, en tres escenarios}
\label{tab:cobertura}
\begin{tabular}{p{7.4cm}rrr}
\hline
Situación del dato & Todos & Sin emisor variable & Estricto \\
\hline
\multicolumn{4}{l}{\textit{Categoría (a): cobertura en una base de la Gobernación}} \\
\quad Total de requisitos & 56 & 23 & 2 \\
\quad La base contiene el dato completo & 1 & 1 & 1 \\
\quad La base contiene una parte del dato & 43 & 19 & 0 \\
\quad Ninguna base declarada contiene el dato & 12 & 3 & 1 \\
\multicolumn{4}{l}{\textit{Categoría (b): mecanismo con la entidad que expide el documento}} \\
\quad Total de requisitos & 167 & 115 & 115 \\
\quad Servicio web declarado, dato completo & 3 & 3 & 3 \\
\quad Servicio web declarado, dato no completo & 36 & 19 & 19 \\
\quad Sin servicio web; consulta pública del dato completo & 2 & 2 & 2 \\
\quad Sin servicio web; consulta pública de vigencia (sin confirmar) & 30 & 30 & 30 \\
\quad Mecanismo no declarado & 96 & 61 & 61 \\
\hline
\end{tabular}

\smallskip
\footnotesize\textit{Nota.} Un requisito es un documento exigido en un trámite. «Sin emisor variable» excluye los documentos que pueden provenir de distintas autoridades; «estricto» excluye además los actos previos de la misma competencia del trámite, que solo aplican a la categoría (a). La consulta pública en línea es una consulta de documento a documento en el portal de la entidad, no un intercambio entre sistemas. Elaboración propia con base en \textcite{gobernaciondelhuila2026datosmaestros}, \textcite{gobernaciondelhuila2026peti}, \textcite{funcionpublica2026visorsuit} y los portales de la Procuraduría, la Contraloría, la Registraduría, el RUES y el RETHUS, consultados el 31 de octubre de 2026.
\end{table}

\subsubsection{Trámites que dependen de información de otras entidades}

De los 85 trámites, 52 exigen al menos un documento de la categoría (b), es decir, dependen de información de otra autoridad. En 48 de ellos al menos uno de esos documentos procede de una entidad con la que no se declara servicio web, y solo en 4 todos proceden de entidades con servicio web declarado. Estos 52 trámites son el insumo para identificar las cadenas de trámites con otras entidades cuya interoperabilidad la Ley 2052 de 2020 obliga a priorizar \parencite[art.~8]{colombia2020ley2052}; el estudio no las desglosó por entidad.

\subsubsection{Estandarización de los documentos en el SUIT}

El SUIT clasifica cada documento como estandarizado o propio del trámite. De los 56 requisitos de la categoría (a), 39 están registrados como propios del trámite y 17 como estandarizados; de los 167 de la categoría (b), 84 están registrados como propios, 81 como estandarizados y 2 tienen ambas clasificaciones, porque el mismo documento se registró con nombres distintos \parencite{funcionpublica2026visorsuit}. Que la mayoría de los documentos de la Gobernación figure como propio del trámite indica que no están asociados a un documento tipo del registro nacional; esta clasificación se reporta como indicador de estandarización y no se usó para asignar dominios.

\subsubsection{Dominios de interoperabilidad}

Los dominios se asignaron solo cuando una fuente da evidencia de ello (Tabla~\ref{tab:dominios}); donde las fuentes callan, el resultado es no determinable y no ausencia del problema. El técnico se asignó cuando la fuente declara que la base o el servicio cubren solo una parte del dato: 43 requisitos de (a) y 36 de (b); en 3 de (b) el dato está completo y, en los demás, no es determinable. El organizacional pudo asignarse a 24 requisitos. En la categoría (a), a los 18 certificados de tiempo de servicio y de salarios de los trámites docentes, porque la fuente declara que esas certificaciones se entregan por oficio o correo, y al certificado de habilitación, cuyo dato está completo en una base y aun así se exige. En la categoría (b), a 3 requisitos del RUT, cuyo dato está completo en una base con servicio web declarado, y a los certificados de antecedentes disciplinarios y fiscales, por la consulta pública descrita. En los demás requisitos, las fuentes no permiten saber si falta un responsable o un procedimiento. El semántico se asignó cuando el campo que registra la base no equivale al dato que certifica el documento: 18 requisitos de (a) y 24 de (b). Se observó además que en 8 de los 10 trámites que exigen el certificado de tiempo de servicio, el SUIT registra una verificación de que se cumple el tiempo de servicio requerido; el registro no indica cómo se hace, por lo que no se usó como evidencia. El político-legal no se asigna en la categoría (a), porque el Decreto Ley 019 de 2012 ya prohíbe pedir los documentos que reposan en la entidad \parencite[art.~9]{colombia2012decreto019}, aunque queda por establecer si, en los trámites docentes, la entidad ante la cual se tramita es la Gobernación o el FOMAG. En la categoría (b) no puede determinarse con fuentes públicas, porque depende, entre otras cosas, de si la Gobernación está vinculada al servicio ciudadano digital de interoperabilidad \pendiente{DATO PENDIENTE: respuesta de la Gobernación a la solicitud 2026PQR00059197}.

\begin{table}[htbp]
\centering
\caption{Requisitos de las categorías (a) y (b) con evidencia de cada dominio de interoperabilidad}
\label{tab:dominios}
\begin{tabular}{p{2.6cm}p{7.4cm}cc}
\hline
Dominio & Criterio de asignación & (a) & (b) \\
\hline
Técnico & La fuente declara que la base o el servicio cubren solo una parte del dato & 43 & 36 \\
Organizacional & El dato está al alcance de la entidad y aun así se pide el documento: entrega manual declarada, base con el dato completo o consulta pública del dato completo & 19 & 5 \\
Semántico & El campo de la base no equivale al dato certificado & 18 & 24 \\
Político-legal & No se asigna en (a); no determinable en (b) & n.\,a. & n.\,d. \\
\hline
\end{tabular}

\smallskip
\footnotesize\textit{Nota.} Categoría (a): \textit{n} = 56; categoría (b): \textit{n} = 167. Un requisito puede tener evidencia de más de un dominio. Los requisitos sin evidencia de un dominio no se cuentan en esa fila y su resultado es no determinable; la ausencia de evidencia no significa que el dominio no intervenga. En la categoría (a), el dominio técnico se refiere a la base de la propia entidad; en la (b), al servicio con la entidad que expide el documento. n.\,a. = no aplica, porque la norma ya prohíbe exigir esos documentos; n.\,d. = no determinable con fuentes públicas. Elaboración propia con base en \textcite{gobernaciondelhuila2026datosmaestros}, \textcite{funcionpublica2026visorsuit} y los portales de la Procuraduría y la Contraloría.
\end{table}

\subsubsection{Documentos que pasan a la Fase 3}

Los criterios de la Fase 3 ordenan primero los documentos de la categoría (a), después por el número de trámites que los exigen y, por último, por la disponibilidad del dato en una base declarada. La Tabla~\ref{tab:candidatos} los aplica a los 16 documentos de la categoría (a) y a los 6 documentos de la categoría (b) que se exigen en 7 o más trámites, el mismo umbral de la Tabla~\ref{tab:repetidos}. En la categoría (a), los dos primeros documentos son los certificados de tiempo de servicio y de salarios, ambos de emisor variable; los documentos que no dependen de ningún supuesto son el certificado de habilitación y la resolución de autorización de etiquetas, que se exigen cada uno en un solo trámite. La lista es un insumo para el diseño y no una selección definitiva.

\begin{table}[htbp]
\centering
\caption{Documentos ordenados según los criterios de la Fase 3}
\label{tab:candidatos}
\small
\begin{tabular}{rp{5.6cm}ccp{4.2cm}}
\hline
\# & Documento & Trámites & Supuesto & Situación del dato \\
\hline
\multicolumn{5}{l}{\textit{Categoría (a)}} \\
1 & Certificado de tiempo de servicio & 10 & EV & Base parcial \\
2 & Certificado de salarios & 8 & EV & Base parcial \\
3 & Resolución de pensión del docente & 8 & MC & Base parcial \\
4 & Acto de retiro del servicio docente & 7 & MC & Base parcial \\
5 & Licencia de funcionamiento & 4 & MC & Base parcial \\
6 & Tornaguía & 4 & EV & No declarada \\
7 & Reportes anuales de cesantías & 3 & EV & No declarada \\
8 & Acta de visita sanitaria & 2 & EV & Base parcial \\
9 & Paz y salvo de la pagaduría del empleador & 2 & EV & Base parcial \\
10 & Resolución de cesantías pagadas & 2 & MC & No declarada \\
11 & Certificado de habilitación & 1 & -- & Base completa \\
12 & Constancia de no sanción por comercio ilegal de medicamentos & 1 & EV & Base parcial \\
13 & Registro profesional en salud & 1 & EV & Base parcial \\
14 & Acta de comité de exoneración del servicio social obligatorio & 1 & EV & No declarada \\
15 & Acta del accidente de trabajo & 1 & EV & No declarada \\
16 & Resolución de autorización de etiquetas & 1 & -- & No declarada \\
\multicolumn{5}{l}{\textit{Categoría (b), documentos exigidos en 7 o más trámites}} \\
17 & Documento de identidad & 30 & -- & Consulta de vigencia (sin confirmar) \\
18 & Certificado de existencia y representación legal & 11 & -- & Servicio web, dato no completo \\
19 & Diploma o acta de grado & 10 & EV & Mecanismo no declarado \\
20 & Registro civil de nacimiento & 9 & -- & Mecanismo no declarado \\
21 & Certificado de administradora de pensiones & 7 & EV & Servicio web, dato no completo \\
22 & Tarjeta profesional & 7 & EV & Servicio web, dato no completo \\
\hline
\end{tabular}

\smallskip
\footnotesize\textit{Nota.} EV = el documento puede provenir de distintas autoridades (emisor variable); MC = acto previo de la misma competencia del trámite; ambos supuestos se explican en la Fase 1. Los documentos de igual número de trámites se ordenan por la disponibilidad del dato. Elaboración propia con base en \textcite{gobernaciondelhuila2026datosmaestros}, \textcite{funcionpublica2026visorsuit} y los portales de las entidades emisoras.
\end{table}
